%%%%%%%%%%%%%%%%%%%%%%%%%%%%%%%%%%%%%%%%%%%%%%%%%%%%%
%			HLAVIČKA								%
%%%%%%%%%%%%%%%%%%%%%%%%%%%%%%%%%%%%%%%%%%%%%%%%%%%%%
\documentclass[openany,12pt]{memoir}

\usepackage[utf8]{inputenc} 
\usepackage[czech]{babel}
\usepackage[T1]{fontenc}
\usepackage[top=1.5cm, bottom=2cm, left=2cm, right=2cm]{geometry}  % --> NASTAVENÍ OKRAJŮ
\usepackage{fancyhdr}
\usepackage{graphicx}
%\usepackage{lmodern}
\usepackage{tgschola}  % --- Silnější font
\usepackage{xwatermark}
\usepackage{xcolor}
\usepackage{changepage}
\usepackage{newunicodechar}



%%%%%% Package na zpěvník
\usepackage[full]{leadsheets}%http://mirrors.nic.cz/tex-archive/macros/latex/contrib/leadsheets/leadsheets_en.pdf   --> dokumentace	
\definesongtitletemplate{empty}{} 
\setchords{
format = \bfseries,   %tučné akordy
minor = {mi},% 
input-notation = {german},%
output-notation = {german}%
}
\definesongtitletemplate{empty}{} 

\newlength{\drop}
\newwatermark[allpages,color=red!50,angle=0,scale=2, xpos=0,ypos=0]{\includegraphics[width=5cm]{obr/pozadi2.jpg}} %--> dvojka na pozadí


%%%% Vlastní příkazy
\newcounter{Slokočet}   %Automatické číslování slok
\newcommand{\mezera}{\vspace*{0.5cm}}   %Horizontální odsazení slok
\newcommand{\stred}{5.2cm}   %%% Na zarovnání slok doprostřed, pozn. automatičtější zarovnávání na střed nejde
\newcommand{\refren}{\mezera \noindent \textbf{R:} } %refrén
\newcommand{\sloka}{\mezera \noindent \addtocounter{Slokočet}{1} \arabic{Slokočet}. } 	%sloka, která se automaticky čísluje
\newcommand{\ssloka}{\mezera \noindent}  % vlastní číslo sloky
\newcommand{\carka}{,\:}
\newcommand{\predehra}{\mezera \hspace*{-1.85cm} \textbf{Předehra:} }
\newcommand{\m}[1]{\color{white}{#1}}  %Pro akordy
\newcommand{\ap}{'} %Pro apostrof
\newcommand{\elipsa}{\kern\fontdimen3\font} %Příkaz pro lepší zacházení s výpustkami (=...); je to vpodstatě jen mezera mezi tečkama výpustky

%%%%%%%%%%%%%%%%%%%%%%%%%%%%%%%%%%%%%%%%%%%%%%%%%%%%%%
%			NÁVOD									 %
%%%%%%%%%%%%%%%%%%%%%%%%%%%%%%%%%%%%%%%%%%%%%%%%%%%%%%
%1. Věci v hlavičce IGNOROVAT
%2. Píseň psát do prostoru mezi \begin{song} a \end{song}
%3. další řádek se značí dvěma odsazeníma (= dvakrát stisknout enter)
%4. \refren vždy na začátku refrenu a \sloka na začátek sloky (automaticky se čísluje)
%  \ = alt gr + q ; [/] = alt gr f/g ; {/} = alt gr + b/n; ^ = alt gr + 3 + mezera
%Cokoliv napíšete do ^{  } se bude brát jako akord
%když se toto bude dotýkat nějakého slova (nebude mezi tím a slovem mezera)
%tak se akordy zjeví nad slovem, ale pište to před slovo
%Když se to nedotýká slova, tak akord lítá ve vzduchu a vytiskne se větší mezera
%První možnost je asi preferovanější
%5. Akordy stačí psát jen do první sloky, když se nezmění -- kytaristi to zvládnou



%\usepackage{subfiles}
%</++>

\begin{document}

\pagestyle{empty}
%\addtocontents{toc}{\protect\thispagestyle{empty}}
%\tableofcontents \thispagestyle{empty}\newpage
\newgeometry{top=1.5cm, bottom = 0cm, left = 2cm, right = 2cm}
	
\newpage

%Pro kompilaci je důležitý, 	aby ten input a \end{document} zůstali na posledních
%dvou řádcích!!!!


\begin{song}{title=\centering Mám doma kočku\\\normalsize Pokáč \vspace*{-0.3cm}}  %% sem se napíše jméno songu a autor
\moveright 0.5cm \vbox{      %Varianta č. 1  ---> Jeden sloupec zarovnaný na střed	

\begin{minipage}[t]{0.48\textwidth}\setlength{\parindent}{0.45cm}  %Varianta č. 2 --> Dva sloupce

Capo 1


\sloka
   ^{G{\color{white}\_\_\_\_}}Jednou mi drahá říká:

   ^{Emi{\color{white}\_\_}}Miláčku, já chci kočku

   ^{C\,\,}Já na to, že chci spíš psa
   
   ^{D}A jeli jsme pro kočku
  
   
\sloka
   V útulku v Měcholupech
   
   Maj dobrej výběr, zdá se
   
   Kočky jsou zcela zdarma
   
   Jen očkování za patnáct set
 
   
\sloka   
   Po dlouhým vybírání
   
   Diskuzích a dvou hádkách
   
   Jedno z koťátek jemně sevře
   
   Můj prst ve svých drápkách

   
\sloka
   Tak volba byla jasná
   
   Má drahá slzy suší
   
   Co si však vezem domů
   
   To ani jeden z nás netuší
 
   
\refren
   Mám doma kočku, mám ji rád
   
   Je to skvělý kamarád
   
   Jen od tý doby, co ji znám
   
   Mám ruce nohy samej šrám
   
   Dávno už se nemazlí
   
   Žít s ní je pro zdraví fakt zlý
   
   Každý den projdu terorem
   
   Mám nejspíš kočku s errorem
 
 
 \sloka
   Po pár dnech soužití
   
   Vše vypadá nevinně
   
   Koťátko v náručí spí
   
   A pak zkadí půlku kuchyně
   
   
\sloka
   Prostě nic, co nečekáš
   
   Když útulek zaručí
   
   Že kočička se za pár týdnů
   
   Na záchůdek naučí
  
  
\end{minipage}\begin{minipage}[t]{0.5\textwidth}\setlength{\parindent}{0.45cm}\vspace*{0.44cm}  % V případě varianty č.2 jde odsud text do pravé části  

\sloka   
   Jenže pár měsíců je pryč
   
   Záchod zeje prázdnotou
   
   Bytem čpí příšernej smrad
   
   Za kočkou lítám s lopatou
 
 
\sloka
   Když jí zkusím pomoct sám
   
   Na záchod ji přemístit
   
   Zvládne mi dlaně rozdrásat
   
   A ještě potřísnit
  
   
\refren
   
   
\sloka
   Za noc nejmíň pětkrát
   
   Jí okno otevřít jdu
   
   Pak práci šéf se zeptá
   
   Proč vždy tak ospalý přijdu
   
   
\sloka
   Že občas hrdě nám
   
   Donese zakouslou myš
   
   Poznáš až ráno když se
   
   S torzem myši v ústech vzbouzíš
  
   
\sloka
   Tak jsme si dobrovolně
   
   Pořídili katastrofu
   
   Stálo nás to patnáct stovek
   
   A jednu zbrusu novou sofu
   
   
\sloka
   Kouše, škrábe, drásá
   
   No prostě je to krása
   
   Tváří se tak roztomile
   
   Zvlášť, když se mi rejpe v žíle
   
   
\refren
   
   
\sloka
   Ráda dlouze sedá v okně
   
   Anebo ve dveřích
   
   Bůh ví, proč nejradši v zimě
   
   Když venku padá sníh
   
   
\sloka
   Když pak v bytě máme mínus dvacet
   
   Tak nás zahřívá
   
   Aspoň pocit, že v tu chvíli naši
   
   krev neprolívá
   
   OU!
   
   
\refren \dots \, MŇAU
   
\end{minipage}   
}
\setcounter{Slokočet}{0}
\end{song}

\end{document}
